% ============================================================================
%  SECTION 1b — Mechanical stimuli: Laplace law & wall shear stress (EARLY)
% ============================================================================
\begin{frame}[t]{An Idealized Blood Vessel}
  A vessel is a \textbf{pressurised, perfused tube}. Two loads act on the wall:
  \textbf{blood pressure} $P$ pushes it outward, and \textbf{blood flow} $Q$ drags
  its inner surface. Cells sense both.
  \vfill
  \begin{center}\idealvessel[0.60]\end{center}
  \vfill
  \begin{center}\small
  $P$ lumen pressure $\to$ \textbf{circumferential stress} $\sigma_\theta$
  \qquad
  $Q$ volumetric flow $\to$ \textbf{wall shear stress} $\tau_w$
  \qquad
  $a$ inner radius, \ $h$ wall thickness
  \end{center}
  \srccite{Humphrey, Cardiovascular Solid Mechanics (2002)}
\end{frame}

\begin{frame}[t]{Two Loads, Two Formulae, Two Adaptations}
  Pressure $P$ sets the \textbf{hoop stress}; flow $Q$ sets the \textbf{wall shear
  stress}. Each has a set-point; rearranging shows the geometry that restores it.
  \vfill
  \begin{columns}[T]
    \begin{column}{0.5\textwidth}
      \begin{center}\textbf{\color{YaleBlue} Pressure} (Laplace)\end{center}
      \slboxeq{2.4em}{eq:laplace}{\sigma_\theta = P\,\frac{a}{h}}
      Hold $\sigma_\theta=\sigma_{\theta,h}$: $\;h=Pa/\sigma_{\theta,h}\propto P$
      --- higher pressure $\Rightarrow$ \textbf{thicker wall}.
    \end{column}
    \begin{column}{0.5\textwidth}
      \begin{center}\textbf{\color{YaleBlue} Flow} (Poiseuille)\end{center}
      \slboxeq{2.4em}{eq:wss}{\tau_w = Q\,\frac{4\mu}{\pi\,a^3}}
      Hold $\tau_w=\tau_{w,h}$: $\;a=(4\mu Q/\pi\tau_{w,h})^{1/3}\propto Q^{1/3}$
      --- higher flow $\Rightarrow$ \textbf{larger lumen}.
    \end{column}
  \end{columns}
  \vfill
  \srccite{Laplace \& Poiseuille; Humphrey, Cardiovascular Solid Mechanics (2002)}
\end{frame}

% ---- interactive analytical exercises (pen & paper) ------------------------
\begin{frame}[t]{}
  \begin{yourturnbox}{Pressure adaptation}
    \qq{Pressure \textbf{doubles} ($P'=2P$), radius $a$ fixed. By how much must the
    wall thicken to restore $\sigma_{\theta,h}$?}
    \qhint{Set $P'a/h'=\sigma_{\theta,h}=Pa/h$; solve for $h'$.}
  \end{yourturnbox}
\end{frame}

\begin{frame}[t]{}
  \begin{answerbox}{Pressure adaptation}
    $P'a/h'=Pa/h$ with $P'=2P$ gives
    \slboxeq{2.2em}{eq:ans:pressure}{h' = 2\,h.}
    \par\smallskip
    {\color{ABlue}\bfseries Reading.} Wall thickens \emph{in proportion} to
    pressure (Grossman) --- radius unchanged, stress restored.
  \end{answerbox}
  \srccite{Grossman, Jones, McLaurin, J Clin Invest (1975)}
\end{frame}

\begin{frame}[t]{}
  \begin{yourturnbox}{Flow adaptation}
    \qq{Flow rises \textbf{8-fold} ($Q'=8Q$). By how much must the radius change to
    restore $\tau_{w,h}$?}
    \qhint{$\tau_w\propto Q/a^3$; set $\tau_w(Q',a')=\tau_{w,h}$ and solve for $a'$.}
  \end{yourturnbox}
\end{frame}

\begin{frame}[t]{}
  \begin{answerbox}{Flow adaptation}
    $(a'/a)^3 = Q'/Q = 8$, so
    \slboxeq{2.2em}{eq:ans:flow}{a' = 8^{1/3}\,a = 2\,a.}
    \par\smallskip
    {\color{ABlue}\bfseries Reading.} Radius grows as the \emph{cube root} of flow
    --- flow-induced outward remodeling; an 8$\times$ flow only doubles the lumen.
  \end{answerbox}
  \srccite{Kamiya \& Togawa, Am J Physiol (1980)}
\end{frame}
